% Drawing 138/07-13 -- folder 138, batch 07, archive page 131 (PDF page 132).
% Transcription: un cercle (la sphère $S$) portant le point $\infty$ et le segment de $0$ à $1$.
% Kind: dessin d'enfant on the sphere (segment 0 - 1). Method 3 (encoded, rendered from the data). Status: ok.
% Notes: A circle (the sphere S) carrying the point infinity and the segment from 0 to 1. Encoded as the trivial dessin: one edge 0 (black) -- 1 (white), one face infinity, genus 0.
% Coordinates: 1 unit = 1 mm = 1/100 of the crop width (bbox in page fractions: [0.136, 0.082, 0.276, 0.17]).
% Redrawn by Claude (claude-opus-5-5), issue #73, 2026-11-07; not checked by a human; not compiled (no TeX in the session).
% Rights: reproduces part of the work -- behind RENDER_TRANSCRIPTS, never published by itself.
\documentclass[tikz,border=3pt]{standalone}
\usepackage{amsmath,amssymb}
\usetikzlibrary{patterns,arrows.meta}
\tikzset{pencil/.style={black!45}, struck/.style={black!35}, mark/.style={circle,draw,inner sep=1.3pt}, pt/.style={circle,fill,inner sep=0.9pt}, opt/.style={circle,draw,fill=white,inner sep=1.1pt}}
\begin{document}
\begin{tikzpicture}[x=1mm,y=1mm,line cap=round,line join=round,font=\scriptsize]
\draw (48.7,-50) ellipse[x radius=35.9, y radius=42.3];
\node[pt] at (48.7,-39.7) {};
\node[inner sep=1pt] at (52.6,-33.3) {\scriptsize $\infty$};
\draw[very thick] plot[smooth, tension=0.55] coordinates {(30.8,-51.3) (44.9,-55.8) (64.1,-52.6)};
\node[pt] at (30.8,-51.3) {};
\node[pt] at (64.1,-52.6) {};
\node[inner sep=1pt] at (28.2,-59) {\scriptsize $0$};
\node[inner sep=1pt] at (69.2,-59) {\scriptsize $1$};
\end{tikzpicture}
\end{document}
